\chapter{Хэрэгжүүлэлт}

\section{Програм хангамжийн хөгжүүлэлт}
Системийн хэрэгжүүлэлтийг Android стандарт болон цэвэр кодын зарчимд тулгуурлан гүйцэтгэв. Үндсэн интерфейс болон үйлдлүүдийн хэрэгжилтийг дараах хэсгүүдэд хуваан харуулав.

\subsection{Үндсэн интерфейс}
Хэрэглэгчийн интерфейсийг Material Design 3 стандартын дагуу оновчилж, нэг гараар ашиглахад тохиромжтой доод навигацийн цэс болон харилцан ярианы нэгдсэн харагдацыг хэрэгжүүлсэн.

\begin{figure}[H]
    \centering
    \includegraphics[width=0.4\textwidth]{Chart/unfolder_main_screen.png}
    \caption{Аппликейшны үндсэн дэлгэц}
    \label{fig:main_screen}
\end{figure}

\begin{figure}[H]
    \centering
    \begin{minipage}{0.45\textwidth}
        \centering
        \includegraphics[width=0.9\textwidth]{Chart/bottom_bar.jpg}
        \caption{Доод навигацийн цэс}
    \end{minipage}\hfill
    \begin{minipage}{0.45\textwidth}
        \centering
        \includegraphics[width=0.9\textwidth]{Chart/all_message_view.jpg}
        \caption{Нийт харилцан ярианы жагсаалт}
    \end{minipage}
\end{figure}

\subsection{Ангилал ба тогтмол илэрхийллийн хэрэгжилт}
Зурвасуудыг автоматаар ангилахын тулд тогтмол илэрхийлэл болон түлхүүр үгсийг удирдах модулийг хэрэгжүүлсэн.

\begin{figure}[H]
    \centering
    \begin{minipage}{0.45\textwidth}
        \centering
        \includegraphics[width=0.9\textwidth]{Chart/edit_category_panel.jpg}
        \caption{Ангилал засах самбар}
    \end{minipage}\hfill
    \begin{minipage}{0.45\textwidth}
        \centering
        \includegraphics[width=0.9\textwidth]{Chart/edit_category_view.jpg}
        \caption{Ангиллын тохиргоо}
    \end{minipage}
\end{figure}

\begin{figure}[H]
    \centering
    \includegraphics[width=0.4\textwidth]{Chart/filter_tag_panel.jpg}
    \caption{Шүүлтүүр болон шошгоны удирдлагын хэсэг}
    \label{fig:filter_panel}
\end{figure}

\subsection{Хадгалсан харагдац буюу Хавтаслах систем}
Хэрэглэгч өөрийн хүссэнээр хавтас үүсгэж, харилцан яриаг тусгаарлах боломжийг бүрдүүлсэн. Хавтас бүрт өнгө оноож, өөрийн гэсэн харагдац (Saved View) үүсгэх боломжтой.

\begin{figure}[H]
    \centering
    \begin{minipage}{0.45\textwidth}
        \centering
        \includegraphics[width=0.9\textwidth]{Chart/create_folder_panel.jpg}
        \caption{Шинэ хавтас үүсгэх}
    \end{minipage}\hfill
    \begin{minipage}{0.45\textwidth}
        \centering
        \includegraphics[width=0.9\textwidth]{Chart/edit_folder_panel.jpg}
        \caption{Хавтас засах самбар}
    \end{minipage}
\end{figure}

\subsection{Хэрэглэгчийн туршлага ба шудрах үйлдэл}
Шудрах үйлдлийг хэрэглэгч өөрийн хэрэгцээнд тааруулан тохируулах боломжтой бөгөөд харилцан ярианы доторх харагдацыг илүү цэгцтэй болгосон.

\begin{figure}[H]
    \centering
    \begin{minipage}{0.45\textwidth}
        \centering
        \includegraphics[width=0.9\textwidth]{Chart/swipe_choose.jpg}
        \caption{Шудрах үйлдэл сонгох}
    \end{minipage}\hfill
    \begin{minipage}{0.45\textwidth}
        \centering
        \includegraphics[width=0.9\textwidth]{Chart/single_conversation_view.jpg}
        \caption{Харилцан ярианы дэлгэрэнгүй харагдац}
    \end{minipage}
\end{figure}

\begin{figure}[H]
    \centering
    \begin{minipage}{0.45\textwidth}
        \centering
        \includegraphics[width=0.9\textwidth]{Chart/swipe_left.jpg}
        \caption{Зүүн тийш шудрах үйлдэл}
    \end{minipage}\hfill
    \begin{minipage}{0.45\textwidth}
        \centering
        \includegraphics[width=0.9\textwidth]{Chart/swipe_right.jpg}
        \caption{Баруун тийш шудрах үйлдэл}
    \end{minipage}
\end{figure}

\section{Хэрэгжүүлэлтийн онцлог}
Төслийн хүрээнд Jetpack Compose ашиглан динамик интерфейсүүдийг хэрэгжүүлсэн нь системийн уян хатан байдлыг нэмэгдүүлсэн. Мөн өгөгдлийн сангийн шилжилт болон нугалдаг төхөөрөмжийн тасралтгүй байдлыг кодын түвшинд бүрэн шийдсэн.
